%& -shell-escape
% !TeX root = thesis-abstract.tex
% !TeX encoding = UTF-8
% !TeX program = XeLaTeX

\documentclass{.local/lib/classes/ndvr-super-article-standalone}

\begin{document}

    Пусть даны два видео — $V(t)$ и $W(t)$.
    Считаем, что видео является частным случаем временного ряда.
    Продолжительность  каждого видео равна $T_{V}$ и $T_{W}$ единиц времени.
    Единицы измерения времени могут не совпадать.
    В каждый момент времени каждый ряд обладает набором характеристик,
    $v_t = V(t_{V})$ и~$w_t = W(t_{W})$  соответственно.
    Размерности $v_t$ и $w_t$  в общем случае неизвестны.
    Из точек временного ряда $V(t)$, в которых наблюдались аномалии, 
    построим временной ряд $(ev_{i})^{k}_{i=1}$.
    \[
        \left\lbrace
            \begin{array}{lcl}
                (ev_{i})^{k}_{i=1} & = &  ev_1, ev_2, \ldots, ev_k; \\
                \left\lbrace ev_1, ev_2, \ldots, ev_k  \right\rbrace & \subset & V 
            \end{array}
        \right.
    \]
    Для $W(t)$ аналогичным образом построим  $(ew_{i})^{n}_{i=1}$.
    \begin{definition}
    $V$ и $W$ выражают одинаковую
    последовательность явлений, если существует $\tser{x}{m}$, 
    такой что $ \tset{x}{m} \subset \tset{ev}{k} $ 
    и~$ \tset{x}{m} \subset \tset{ew}{n}$.
    \[
        V \esim  W \Longleftrightarrow  
        \exists \tser{x}{m}; \quad \left\lbrace
            \begin{array}{lclll}
                \tset{x}{m} & \subset 
                    & \tset{ev}{k} 
                    & \subset & V; \\
                \tset{x}{m} & \subset 
                    & \tset{ew}{n} 
                    & \subset & W 
            \end{array}
        \right.
    \]
    \end{definition}
    \begin{definition}
    $V$ и $W$ являются нечеткими дубликатами друг друга,
    если существует $\tser{x}{m}$, 
    такой что $ \tset{x}{m} \subset \tset{v}{k} $ 
    и~$ \tset{x}{m} \subset \tset{w}{n}$.
    \[
        V \ndsim  W \Longleftrightarrow 
        \exists \tser{x}{m}; \quad \left\lbrace
            \begin{array}{lclll}
                \tset{x}{m} & \subset 
                    & \tset{v}{k} 
                    & \subset & V; \\
                \tset{x}{m} & \subset 
                    & \tset{w}{n} 
                    & \subset & W 
            \end{array}
        \right.
    \]
    \end{definition}
    \begin{statement}[О дубликатах]
        Нечеткие дубликаты видео выражают 
        одинаковую последовательность явлений.
        \[
             V \esim  W \Longrightarrow  V \ndsim  W
        \]
    \end{statement}

\end{document}
 
